\section{Related Work}
\textbf{Quantization-Aware Training (QAT)} is often used to obtain quantized models that are adapted in downstream tasks \citep{peri2020deploying, liu2023llm}. It involves quantization and full model fine-tuning at the same time. However, QAT requires massive training cost, such as the gradient and optimization state. Moreover, it is difficult to compute the gradient of quantized weights.
Our method, with the help of LoRA, sidesteps the aforementioned issues, providing a light approach for downstream task adaptation.

\vskip2pt

\noindent \textbf{Post-Training Quantization (PTQ)} is a category of popular quantization frameworks \citep{frantar2022gptq, xiao2023smoothquant}, which can also be used for task adaptation. It calibrates the high-precision model with a small subset of the training dataset. Therefore, the subsequent quantization is guided by the training dataset, providing task-specific quantized models. Besides, it does not involve any gradient backpropagation, so it is cost-efficient. However, it usually results in lower accuracy compared to QAT. 

\section{Conclusion}
We propose {\OurAlg}, a quantization framework for LLMs, which alternatively applies quantization and low-rank approximation to the original high-precision pre-trained weights, to obtain an initialization for the subsequent LoRA fine-tuning. Experiments on natural language understanding, question answering, summarization, and natural language generation show that our framework remarkably surpasses existing methods, e.g., QLoRA, for quantizing encoder-only, encoder-decoder, and decoder-only models. We have not observed our method exhibiting worse performance over QLoRA. Moreover, our quantization framework demonstrates effectiveness and robustness particularly in low-bit quantization regimes, e.g., the 2-bit level.
